\chapter{Sclerotherapy}

\section*{Overview}
Sclerotherapy is a procedure that treats varicose and spider veins by injecting a solution that causes the vein to collapse and fade.

\section*{Alternative Names}
Varicose vein injection therapy; sclerosing injection; spider vein treatment

\section*{Principle}
A sclerosing agent is injected into a vein, irritating its lining and causing the vein to close and scar, which redirects blood to healthier veins. The treated vein gradually fades and is reabsorbed.

\section*{History}
Sclerotherapy has been practiced in various forms since the 19th century and remains a first-line treatment for small varicose and spider veins.

\section*{Why It Is Done}
Performed for symptomatic spider veins and small varicose veins causing aching, burning, or swelling, and for cosmetic improvement.

\section*{Is This a Primary or Secondary Test or Procedure}
Primary treatment for spider veins and small varicose veins.

\section*{How It Is Performed}
The vein is injected with a liquid or foam sclerosant. Foam sclerosant is often used for larger veins. Compression is applied afterward. Several sessions may be needed.

\section*{Preparation}
A venous ultrasound may precede treatment; sun exposure, certain medications, and compression hosiery recommendations are reviewed.

\section*{Contraindications and Precautions}
Deep vein thrombosis, severe peripheral arterial disease, allergy to the sclerosant, and pregnancy are contraindications.

\section*{Risks}
Pain, bruising, skin discoloration, blood clots, ulceration, and rarely allergic reactions.

\section*{Aftercare and Recovery}
Compression stockings are worn for days to weeks; walking is encouraged, and treated veins fade over several weeks.

\section*{Outcome and Follow-up}
Improvement is judged by fading of veins and symptom relief; persistent vessels may require repeat injections.

\section*{Expected Results}
Fading or disappearance of treated veins with no complications.

\section*{Follow-up}
Follow-up visits and ultrasound if large truncal veins are present.

\section*{When to Seek Medical Care}
Follow the procedure team's specific instructions. Seek urgent medical assessment for severe or worsening pain, uncontrolled bleeding, fever or signs of infection, trouble breathing, chest pain, fainting, new weakness or confusion, inability to urinate, or any rapidly worsening symptom. The appropriate warning signs depend on the procedure and the patient's medical condition.

\section*{References}
\begin{enumerate}
  \item MedlinePlus. Sclerotherapy. U.S. National Library of Medicine; 2025.
  \item Current Medical Diagnosis and Treatment. McGraw-Hill; 2025.
\end{enumerate}

\clearpage
